\documentclass[10pt]{article}
\usepackage[T1]{fontenc}
\usepackage[utf8]{inputenc}
\usepackage[french]{babel}
\usepackage{amsmath,amssymb,mathrsfs,mathpazo}
\usepackage{geometry}
\geometry{a4paper,margin=2cm}
\usepackage{graphicx,array,tabularx,tabularray,colortbl,multirow,multicol,booktabs,enumitem,pifont,ulem,fancybox,diagbox,cancel}
\usepackage{pstricks,pstricks-add,pst-plot,pst-tree,pst-node,pst-text,pst-func,pst-math,pst-bspline,pst-3dplot,pst-eucl,pst-solides3d}
\let\pscalculate\undefined
\usepackage{pst-bezier}
\usepackage{tikz,pgfplots,tkz-tab}
\usetikzlibrary{arrows,arrows.meta,patterns,calc,positioning,trees,decorations.markings,shapes.geometric,shapes.misc,shapes.arrows,intersections,decorations.pathmorphing,decorations.pathreplacing,backgrounds,shadows,babel}
\pgfplotsset{compat=1.18}
\usepgfplotslibrary{fillbetween}
\usepackage[np]{numprint}
\usepackage{listings,makecell,pas-tableur,twoopt}
\let\barre\undefined
\usepackage[pstricks]{tablvar}
\definecolor{darkgreen}{RGB}{0,100,0}
\providecommand{\red}{\color{red}}
\providecommand{\blue}{\color{blue}}
\providecommand{\green}{\color{green}}
\providecommand{\black}{\color{black}}
\providecommand{\white}{\color{white}}
\providecommand{\gray}{\color{gray}}
\providecommand{\cyan}{\color{cyan}}
\providecommand{\magenta}{\color{magenta}}
\providecommand{\yellow}{\color{yellow}}
\providecommand{\violet}{\color{violet}}
\providecommand{\orange}{\color{orange}}
\providecommand{\euro}{\texteuro}
\providecommand{\EUR}{\texteuro}
\providecommand{\ieme}{\textsuperscript{e}}
\providecommand{\degres}{\ensuremath{^\circ}}
\providecommand{\degree}{\ensuremath{^\circ}}
\providecommand{\vv}[1]{\overrightarrow{#1}}
\providecommand{\Coord}[1]{#1}
\providecommand{\decosix}{}
\providecommand{\starredbullet}{\ensuremath{\bullet}}
\pagestyle{empty}
\setlength{\parindent}{0pt}
\usepackage{adjustbox}\DeclareUnicodeCharacter{2236}{:}\DeclareUnicodeCharacter{2190}{\ensuremath{\leftarrow}}
\providecommand{\R}{\mathbb{R}}
\providecommand{\N}{\mathbb{N}}
\providecommand{\Z}{\mathbb{Z}}
\providecommand{\Q}{\mathbb{Q}}
\providecommand{\D}{\mathbb{D}}
\providecommand{\C}{\mathbb{C}}
\providecommand{\e}{{\mathrm{e}}}
\providecommand{\vect}[1]{\overrightarrow{#1}}
\providecommand{\vectt}[1]{\overrightarrow{\mathrm{#1}}}
\providecommand{\vvt}[1]{\overrightarrow{\mathrm{#1}}}
\providecommand{\np}[2][]{{#2~#1}}
\providecommand{\vv}[1]{{\overrightarrow{#1}}}
\renewcommand{\d}{{\mathrm{d}}}
\providecommand{\degres}{{^\circ}}
\providecommand{\degre}{{^\circ}}
\providecommand{\blue}{\color{blue}}
\providecommand{\red}{\color{red}}
\providecommand{\enskip}{\quad}
\providecommand{\llbracket}{\lbrack\!\lbrack}
\providecommand{\rrbracket}{\rbrack\!\rbrack}
\providecommand{\euro}{{\text{€}}}
\providecommand{\no}{{\text{n°}}}
\providecommand{\interval}[2]{{[#1;#2]}}
\providecommand{\intervFO}[2]{{[#1;#2[}}
\providecommand{\dotfill}{\ldots}

\providecommand{\ds}{\displaystyle}
\providecommand{\textding}[1]{#1}
\providecommand{\Oij}{\(\mathrm{O};\vec{\imath},\vec{\jmath}\)}
\providecommand{\Oijk}{\(\mathrm{O};\vec{\imath},\vec{\jmath},\vec{k}\)}
\providecommand{\Ouv}{\(\mathrm{O};\vec{u},\vec{v}\)}
\providecommand{\intervFF}[2]{\ensuremath{[#1;#2]}}
\renewcommand{\euro}{\ifmmode\text{\texteuro}\else\texteuro\fi}
\DeclareUnicodeCharacter{00B0}{\ifmmode^{\circ}\else\textdegree\fi}
\DeclareUnicodeCharacter{20AC}{\ifmmode\text{\texteuro}\else\texteuro\fi}
% Énoncés : Éducation nationale. Transcription et corrigés : APMEP, Denis Vergès et contributeurs. Découpage et classement : CoopMaths, Sébastien Lozano, Éric Elter et contributeurs. Adaptation de la présentation : Maths Explorer.
% CC BY-SA : https://coopmaths.fr/www/about/
% https://coopmaths.fr/alea/static/eam/2026/tex/eam_technologique_2026_06_polynesie_3.tex
\begin{document}
\noindent\textit{APMEP / CoopMaths. CC BY-SA. Présentation adaptée par Maths Explorer.}\par\medskip
On considère la fonction $f$ définie et dérivable sur $[-4 ; 6]$ par
\[f(x) = -0,5x^2 +  x + 7,5.\]

On note $f'$ la fonction dérivée de $f$.

Dans le repère ci-dessous, on a tracé la courbe représentative de la fonction $f$ et sa tangente au point d'abscisse 3 qui passe par les points de coordonnées $(3 ; 6)$, $(4 ; 4)$ et $(6 ; 0)$.

\begin{center}
\begin{tikzpicture}
\begin{axis}[/pgf/number format/.cd, use comma,1000 sep={\,},
				width=13cm,height=10cm,
				xmin=-4.2,xmax=7.2,ymin=-4.5,ymax=8.2,
				xtick={-4,-3,...,7},
				minor x tick num =0,
				ytick={-5,-4,...,8},
				minor y tick num =0,
				tick label style={font=\small},
				grid=both, axis lines=center,
				grid style={gray!35,line width=0.3pt},
				minor grid style={gray!20,line width=0.2pt},
				enlargelimits=false]
	\addplot[color=blue,domain=-4:8,smooth,samples=200,line width=1.2pt] {(-0.5*x^2+ x + 7.5)};
	\addplot[color=red,domain=-4:8,smooth,samples=2,line width=1.2pt] {-2*(x-6)};
	\node[draw,cross out, line width=1.1pt,minimum width=3pt, minimum height = 3pt] at (axis cs:3,6){};
	\node[draw,cross out, line width=1.1pt,minimum width=3pt, minimum height = 3pt] at (axis cs:4,4){};
	\node[draw,cross out, line width=1.1pt,minimum width=3pt, minimum height = 3pt] at (axis cs:6,0){};
	\node[below left] at (axis cs:0,0){\small 0};
\end{axis}
\end{tikzpicture}
\end{center}
\begin{enumerate}
\item Montrer que, pour tout $x$ appartenant à l'intervalle $[-4 ; 6]$,
\[
f(x)=(0,5x + 1,5)(-x+5)
\]
\end{enumerate}

\textbf{\emph{Les justifications demandées pour les trois questions suivantes peuvent se faire par des calculs ou par des lecture graphiques.}}


\begin{enumerate}[resume]
\item Déterminer le nombre dérivé en 3.
\item Pour chacune des affirmations suivantes, indiquer si elle est vraie ou fausse.

Chaque réponse doit être justifiée.

Une réponse non justifiée ne rapporte aucun point.
\begin{enumerate}
\item \textbf{Affirmation a.} : les solutions de l'inéquation $f(x)\leqslant 0$ sur $[-4 ; 6]$ sont les nombres réels appartenant à l'intervalle $[-3 ; 5]$.


\item \textbf{Affirmation b.} : $f'(x)$ est positif sur l'intervalle $[1 ; 4]$.
\end{enumerate}
\end{enumerate}

\end{document}
